\section{理想回路・有限実行予算での精度}\label{sec:finite}
前節までの方式は成功サンプル数を固定し、期待費用を評価する。
ここでは回路・応答を理想的とし、ノイズを導入せず、全フィルタ試行を
$T\geq1$回で打ち切る。棄却時にも読出し1を予約して1試行を
$c_\psi+4n+1$単位と数えれば、総予算$B$から許される回数は
$T_n(B)=\lfloor B/(c_\psi+4n+1)\rfloor$である。
この保守的な計上を全候補に共通に課す。

実装成功重みを$w$、成功確率を$z=\sum_i p_iw(\lambda_i)$、
成功後の平均を$\mu=\sum_i p_i\lambda_iw(\lambda_i)/z$とする。
Gram実験では$w=\kappa^2P^2$。1試行の結果は、棄却・成功かつ0・成功かつ1の
3種類で、確率はそれぞれ$1-z,z(1-\mu),z\mu$である。よって
\[
 M\sim\operatorname{Bin}(T,z),\qquad
 K\mid M=m\sim\operatorname{Bin}(m,\mu).
\]
$M=0$なら値を返さず、$M>0$なら$\widehat e=K/M$を返す。
事前区間へのクリッピングはしない。値を返せない確率と、返せた場合の条件付き
二乗誤差を別々に示し、成功数を総実行数と混同しない。

\begin{proposition}[有限試行の厳密な法則]\label{prop:finite-law}
$q=1-z$、$C_T(z)=1-q^T$とし、
\[
 H_T(z)=\sum_{m=1}^T\binom Tm z^mq^{T-m}/m,\qquad h_T(z)=H_T(z)/C_T(z)
\]
とおく。このとき
\begin{equation}\label{eq:finite-mse}
 \Pr(M=0)=q^T,\qquad
 R_T=\mathbb E[(\widehat e-e)^2\mid M>0]
 = (\mu-e)^2+\mu(1-\mu)h_T(z).
\end{equation}
無出力も失敗と数えた許容誤差$E$での失敗確率は厳密に
\begin{equation}\label{eq:finite-failure}
 q^T+\sum_{m=1}^T\binom Tm z^mq^{T-m}
 \Pr\{\lvert\operatorname{Bin}(m,\mu)/m-e\rvert>E\}
\end{equation}
である。$h_T$は成功確率$z$に関して単調非増加で、
$h_T(1)=1/T$、$z\downarrow0$で$h_T(z)\to1$。
\end{proposition}
\begin{proof}
$M=m>0$で条件づけると、標本平均の平均は$\mu$、分散は
$\mu(1-\mu)/m$。二乗誤差と失敗事象をそれぞれ$M$で平均すれば両式が得られる。
正の二項分布は$\log(z/(1-z))$を自然母数とする指数型分布族なので、有限和を微分して
\[
 h_T'(z)=\frac{\operatorname{Cov}(1/M,M\mid M>0)}{z(1-z)}
 =\frac{1-h_T(z)\mathbb E[M\mid M>0]}{z(1-z)}\leq0.
\]
最後はJensen不等式による。端点の極限は有限和から従う。
\end{proof}

これらの逆二項モーメントと統計的事実は既知であり\cite{Znid,Drakakis}、
ショット雑音や$1/\sqrt T$則を新規性とはしない。
$M$を期待値$Tz$に置き換える近似は有限回では厳密でない。
$Tz$が小さい場合、値が返ったという条件下では成功数がほぼ1になるが、
同時に無出力が多くなる。$T$を増やすと分散項は減る一方、
固定フィルタのバイアス二乗は残る。MSEと固定信頼度の保証は異なる指標であり、
後者は式\eqref{eq:finite-failure}で評価する。

\begin{corollary}[固定締切での一様信頼度]\label{cor:deadline}
全入力で$z\geq s>0$、$0\leq\mu-e\leq b<E$なら、無出力も失敗と数えて
\begin{equation}\label{eq:deadline-confidence}
 \Pr\{\text{無出力 または }|\widehat e-e|>E\}
 \leq\min\{1,\,2[1-s(1-e^{-2(E-b)^2})]^T\}.
\end{equation}
よって信頼度$1-\delta$に十分な確定試行数は
\[
 T\geq\left\lceil\frac{\log(2/\delta)}
 {-\log[1-s(1-e^{-2(E-b)^2})]}\right\rceil.
\]
\end{corollary}
\begin{proof}
$M=m>0$なら$e$から$E$を超える偏差は$\mu$から少なくとも$E-b$の偏差を必要とし、
Hoeffding上界は$2e^{-2m(E-b)^2}$。$m=0$でも右辺は2なので失敗の上界となる。
二項分布の母関数で平均し、$z\geq s$を使えば従う。
\end{proof}
これは既知の集中不等式と母関数の帰結であり、新しい集中不等式とはしない。
バイアス・成功確率の証明書を、期待費用だけでなく固定締切での信頼度に接続する。

\begin{theorem}[条件付きリスクの3固有値への縮約]\label{thm:three}
$T,a,u$を固定し、連続な重み$0\leq w\leq1$が$[0,u]$上で正とする。
許される全ての有限スペクトル・初期重みについての式\eqref{eq:finite-mse}の最大値は
\begin{equation}\label{eq:three}
 a\delta_e+(1-a)\{\theta\delta_x+(1-\theta)\delta_y\},\qquad
 0\leq e\leq u,\quad e\leq x,y\leq1,\quad0\leq\theta\leq1
\end{equation}
で達成される。固有値の重複を許すので、目標質量は$a$を超えてもよい。
\end{theorem}
\begin{proof}
任意の入力から$e$での質量$a$を確保し、残りを元の有限支持上の確率測度$r$とする。
必要なら$r$にも$e$を含める。この入力の事後平均$\mu$を固定し、
\begin{equation}\label{eq:moment-lp}
 \int dr=1,\qquad
 \int(x-\mu)w(x)\,dr(x)=\frac{a}{1-a}(\mu-e)w(e),\qquad r\geq0
\end{equation}
の下で$\int w\,dr$を最小化する。元の測度が実行可能である。
等式制約は高々2本なので、有限LPの最適基本解は高々2つの正の重みで表せる。
新しい入力は$\mu$を保ち$z$を減らすため、$h_T$の単調性から条件付きMSEは減らない。
目標の1点を加えて高々3固有値になる。逆に式\eqref{eq:three}は許される入力である。
そのパラメータ集合はコンパクトで、$z\geq a\min_{[0,u]}w>0$よりリスクは連続。
したがって最大値を達成する。
\end{proof}

これは既知のモーメント集合・基本解の議論\cite{Winkler,Pinelis}を、このリスクへ適用した
正確な低次元化である。バイアス・分散・成功確率の別々の最悪例が同時に成立すると
仮定するものではない。常に3点が必要とは主張しない。また格子性を持つ失敗確率、
クリッピングした推定量、別の停止則に同じ縮約が成り立つとは述べない。

\paragraph{計算可能な一様上界と選択.}
全入力で$z\geq s>0$、$0\leq\mu-e\leq b$なら、
$t_b=\min\{u+b,1/2\}$、$v_b=t_b(1-t_b)$に対し
\begin{equation}\label{eq:finite-upper}
 R_T\leq b^2+v_bh_T(s),\qquad \Pr(M=0)\leq(1-s)^T.
\end{equation}
上端零化の場合、無出力の上界は最小目標透過率と上端の2点で達成する。
MSE上界は各最悪量が両立しない場合があり、一般に等号とはしない。
\[
 A_T=\frac{1-(1-s)^{T+1}}{(T+1)s},\quad
 B_T=\frac{1-[1+(T+1)s](1-s)^{T+1}}{(T+1)(T+2)s^2}
\]
とおくと、厳密な計算用上界は
\begin{equation}\label{eq:inverse-upper}
 h_T(s)\leq\min\left\{1,
 \frac{A_T+3B_T-\frac52(1-s)^T}{1-(1-s)^T}\right\}.
\end{equation}
実際、$m\geq1$で$1/m\leq1/(m+1)+3/[(m+1)(m+2)]$。
右辺2項の全二項分布での期待値が$A_T,B_T$であり、$m=0$分を除くと
$5(1-s)^T/2$を差し引く。実装では有理入力と外向き180bit累乗で包絡する。
数値積分は厳密法則の照合用に限り、選択の証明書には用いない。

有限候補群について$T_n(B)$でこのMSE上界を最小化し、
無出力確率の保証上界を0.05以下に制限する。
量子計算をしない事前中点も、最悪MSE$u^2/4$・無出力0の候補として含める。
最適性は候補群の計算した上界に対してのみであり、一般的な量子推定の最適性ではない。
ポストセレクションの既存議論\cite{Combes}と同様、棄却された実行を除外した
条件付き精度だけから優位性を結論しない。
